% ===== 原文 pp.13–14 =====
%
% tr-02 / chunk b —— 覆盖原文 p.13–p.14；公式 tag 2.50–2.60
% 注：本 chunk 首段起于原文 p.12 末两行、续于 p.13 首行（该段跨页），为保持句子完整整段排在此处。
% ---------------------------------------------------------------------

本综述后面会用到双 Gamma 函数的以下性质。第一，假设 $\epsilon_{1,2}$ 为实数，则当 $x$ 为实数时
$\Gamma_2(x|\epsilon_1,\epsilon_2)$ 为实数。
% [原文 p.13]
然后，由解析延拓，我们有
\begin{equation}
  \Gamma_2(x^*|\epsilon_1,\epsilon_2) = \Gamma_2(x|\epsilon_1,\epsilon_2)^*.
  \tag{2.50}
\end{equation}

我们需要的另一个关系是
\begin{equation}
  \Gamma_2(x+\epsilon_1|\epsilon_1,\epsilon_2)\Gamma_2(x+\epsilon_2|\epsilon_1,\epsilon_2)
  = x\,\Gamma_2(x|\epsilon_1,\epsilon_2)\Gamma_2(x+\epsilon_1+\epsilon_2|\epsilon_1,\epsilon_2).
  \tag{2.51}
\end{equation}

% [Tutor Note: p.13, 式 (2.50)–(2.51)：这是 Barnes 双 Gamma 函数的两条基本性质（实性与移位关系）。
%   注意 (2.51) 两边都是四因子乘积，不能把 $\Gamma_2(x)$ "约掉"——它是 $\Gamma_2$ 的平移递推，
%   相当于 $\Gamma(x+1)=x\Gamma(x)$ 在双 Gamma 情形的推广。把它与 (2.48) 定义的 $\Upsilon$ 联立，
%   即可得到 $\Upsilon$ 函数的平移递推（这是 DOZZ 公式各因子解析结构的来源）。]

在继续之前，请注意两点函数与三点函数在交换 $b\to b^{-1}$ 下本质不变，因为 (2.46) 的第一行
可以被吸收进主场的定义之中。从使用 (2.43) 的拉格朗日描述来看，这种 $b$ 反演下的不变性
相当令人惊讶，因为它在经典层面完全看不到。在人们最初求解 Liouville 理论时，
这一不变性起了关键作用。

% [Tutor Note: p.13, "第一行可以被吸收进主场的定义"：$b\to b^{-1}$ 时
%   $[\pi\mu\gamma(b^2)b^{2-2b^2}]^{(Q-\alpha_1-\alpha_2-\alpha_3)/b}$ 会乘上一个只依赖各 $\alpha_i$
%   的因子（因为 $Q=b+1/b$ 本身在 $b\to b^{-1}$ 下不变），而三点函数允许对每个主场作常数重标度
%   $V_\alpha\to c_\alpha V_\alpha$；把这些因子分配给三个主场，就能把这一行"吸收"掉，剩下的
%   $\Upsilon$ 组合（含真正的 $\alpha_i$ 依赖）保持不变。两点函数同理。]

\subsubsection{四点函数：一般结构}

在二维规范理论的情形，知道圆柱配分函数和三孔球面的配分函数，就足以算出任意曲面上、
任意孔数的配分函数。类似地，对 Liouville 理论而言，知道两点函数与三点函数，
就足以得到任意曲面上、任意多个算符插入的关联函数。

% [Tutor Note: p.13, 这一段是全文的结构性类比：§2.1 用"圆柱配分函数 + 三孔球面配分函数"这两块积木
%   算出任意曲面上的 $q$-形变杨–米尔斯配分函数；本节则用"两点函数 + 三点函数"这两块积木算出任意
%   曲面上的 Liouville 关联函数。两边的积木一一对应：规范理论那边是特征标与缝合，Liouville 这边是
%   DOZZ 三点函数 (2.46) 与 Virasoro 共形块（(2.55) 里插入的完备态集）。§5 的 2d/4d 对应就是把这
%   两套积木逐块对上，读后文时值得回头对照这一段。]

我们通过计算四点函数来说明这一方法。利用共形不变性，我们可以把三个算符放在 $z=\infty$、
$z=1$ 与 $z=0$。设第四个算符的位置为 $q$，我们想要求出
\begin{equation}
  \langle \ee^{2(Q-\alpha_1)\phi(\infty)}\ee^{2\alpha_2\phi(1)}
  \ee^{2\alpha_3\phi(q,\bar q)}\ee^{2\alpha_4\phi(0)}\rangle .
  \tag{2.52}
\end{equation}
做法是在 $\ee^{2\alpha_2\phi(1)}$ 与 $\ee^{2\alpha_3\phi(q,\bar q)}$ 之间插入完备态集，见图 6。

\begin{center}
\begin{minipage}{\textwidth}
\centering
\includegraphics[width=0.5\textwidth]{fig6.png}\par\vspace{2pt}
\figtitle{6}{四穿刺球面由两个三穿刺球面粘合而成}
\end{minipage}
\end{center}

% [Tutor Note: p.13, 图 6 与 (2.52)：左图是粘合前的两个三穿刺球面（虚线处即待粘合的一对穿刺点：
%   左边球面在 $z=0$、右边球面在 $z'=\infty$），右图是粘合后在 $z$ 坐标下的四穿刺球面，穿刺点为
%   $\infty,1,q,0$；管道旁标的 $q$ 就是等同关系 $zw=q$ 的参数，与 (2.60) 里共形块的展开变量是同一个
%   $q$。原文图注里 "two three-punctured sphere" 用的是单数，中译按语义作"两个三穿刺球面"。]

在几何层面，我们有一个带局部坐标 $z$ 的三穿刺球面，以及另一个带局部坐标 $z'$ 的三穿刺球面。
两者的穿刺点都位于 $z,z'=\infty,1,0$。我们把 $z=0$ 与 $z'=\infty$ 连接起来，做法如下。
$z'=\infty$ 处的局部坐标最好改想成 $w=0$，其中 $wz'=1$。于是，把一个位于 $z=0$、另一个位于
$w=0$ 的两个穿刺点以参数 $q$ 粘合，就是作等同 $zw=q$。最终我们得到一个坐标 $z$ 下的
四穿刺球面，穿刺点位于 $z=\infty,1,q$ 与 $0$。

完备态集由下列算符给出
\begin{equation}
  O_{p,\{n\},\{\tilde n\}} := L_{-n_1}L_{-n_2}\cdots L_{-n_k}
  \bar L_{-\tilde n_1}\bar L_{-\tilde n_2}\cdots \bar L_{-\tilde n_{\tilde k}}\,
  \ee^{2(Q/2+\ii p)\phi}
  \tag{2.53}
\end{equation}

% [Tutor Note: p.13, 式 (2.53)：完备态集里 $p$ 是连续参数，求和要配 $\int_0^\infty\dd p$（见 (2.55)），
%   这是非有理共形场理论（连续谱）与有理共形场理论（有限多个表示）的关键区别。另注意 (2.53) 的
%   归一化并非正交归一，它由 (2.56) 的 Gram 矩阵 $G$ 定义，正因如此 (2.55) 里才会出现 $G$ 的逆矩阵。
%   后裔指标是"降序正整数串"，且允许 $k=0$ 或 $\tilde k=0$（不动后裔、只留主场）。]

% [原文 p.14]
其中 $p\ge 0$ 为实数且为正，而正整数
\begin{equation}
  n_1\ge n_2\ge\cdots\ge n_k,\qquad
  \tilde n_1\ge\tilde n_2\ge\cdots\ge\tilde n_{\tilde k}
  \tag{2.54}
\end{equation}
标定 Virasoro 代数的后裔场。注意 $k$ 或 $\tilde k$ 可以为零。于是四点函数分解为
\begin{equation}
  \begin{aligned}
    &\langle \ee^{2(Q-\alpha_1)\phi(\infty)}\ee^{2\alpha_2\phi(1)}
      \ee^{2\alpha_3\phi(z,\bar z)}\ee^{2\alpha_4\phi(0)}\rangle
      = \int_0^\infty \dd p \!\!\sum_{\{n\},\{\tilde n\},\{n'\},\{\tilde n'\}} \\
    &\qquad \times \langle \ee^{2(Q-\alpha_1)\phi(\infty)}\ee^{2\alpha_2\phi(1)}
      O_{p,\{n\},\{\tilde n\}}\rangle\,
      z^{-\sum n}\bar z^{-\sum\tilde n}|z|^{-2\Delta_p} \\
    &\qquad \times G^{\{n\},\{n'\}}G^{\{\tilde n\},\{\tilde n'\}}
      \langle O_{p,\{n'\},\{\tilde n'\}}^{\dagger}\ee^{2\alpha_3\phi(1)}
      \ee^{2\alpha_4\phi(0)}\rangle
  \end{aligned}
  \tag{2.55}
\end{equation}
其中 $\Delta_p=p^2+Q^2/4$，而 $G^{\{n\},\{n'\}}_pG^{\{\tilde n\},\{\tilde n'\}}_p$ 是下式的逆矩阵：
\begin{equation}
  \resizebox{0.985\textwidth}{!}{$\displaystyle
  \begin{aligned}
    G_{p,\{n\},\{n'\}}G_{p,\{\tilde n\},\{\tilde n'\}}
    &= \langle O_{p,\{n'\},\{\tilde n'\}}^{\dagger}O_{p,\{n\},\{\tilde n\}}\rangle \\
    &= \langle \ee^{(Q-2\ii p)\phi(\infty)}L_{n'_1}\cdots L_{n'_{k'}},
      \bar L_{\tilde n'_1}\cdots \bar L_{\tilde n'_{\tilde k'}},
      L_{-n_1}\cdots L_{-n_k}\bar L_{-\tilde n_1}\cdots \bar L_{-\tilde n_{\tilde k}}
      \ee^{(Q+2\ii p)\phi(0)}\rangle .
  \end{aligned}$}
  \tag{2.56}
\end{equation}
我们很快就会看到，$G_{p,\{n\},\{n'\}}$ 可以仅用 Virasoro 代数算出。

我们也可以写
\begin{equation}
  \langle \ee^{2(Q-\alpha_1)\phi(\infty)}\ee^{2\alpha_2\phi(1)}O_{p,\{n\},\{\tilde n\}}\rangle
  = l_{\alpha_1,\alpha_2,\alpha,\{n\}}l_{\alpha_1,\alpha_2,\alpha,\{\tilde n\}}
  C(\alpha_1,\alpha_2,\alpha)
  \tag{2.57}
\end{equation}
\begin{equation}
  \langle O_{p,\{n'\},\{\tilde n'\}}^{\dagger}\ee^{2\alpha_3\phi(1)}
  \ee^{2\alpha_4\phi(0)}\rangle
  = r_{\{n'\},Q-\alpha,\alpha_3,\alpha_4}r_{\{\tilde n'\},Q-\alpha,\alpha_3,\alpha_4}
  C(Q-\alpha,\alpha_3,\alpha_4)
  \tag{2.58}
\end{equation}
其中 $\alpha=Q/2+\ii p$，而函数 $l$ 与 $r$ 同样可以仅用 Virasoro 代数算出。

% [Tutor Note: p.14, 式 (2.57)–(2.58)：这两式把"含后裔的三点函数"拆成一个纯代数因子（$l$ 或 $r$）
%   乘上"主场三点函数" $C$。$l$ 型的因子来自 $z=\infty$ 与 $z=1$ 两点处的两个主场与中间态的后裔，
%   $r$ 型的因子来自 $z=0$ 处的主场；拆完之后，Virasoro 代数的全部信息都集中在 $l,r,G$ 里，
%   而 DOZZ 部分只剩 $C$。下一小节的任务就是把 $l,r,G$ 的前几阶显式算出来。]

把这些关系代回 (2.55)，我们得到
\begin{equation}
  \begin{aligned}
    &\langle \ee^{2(Q-\alpha_1)\phi(\infty)}\ee^{2\alpha_2\phi(1)}
      \ee^{2\alpha_3\phi(q,\bar q)}\ee^{2\alpha_4\phi(0)}\rangle
      = \int_0^\infty \dd p\,|q|^{2\Delta_p} \\
    &\qquad \times C\!\left(\alpha_1,\alpha_2,\frac{Q}{2}+\ii p\right)
      C\!\left(\frac{Q}{2}-\ii p,\alpha_3,\alpha_4\right)
      F_{\alpha_1,\alpha_2,Q/2+\ii p,\alpha_3,\alpha_4}(q)
      \overline{F_{\alpha_1,\alpha_2,Q/2+\ii p,\alpha_3,\alpha_4}(q)}
  \end{aligned}
  \tag{2.59}
\end{equation}
其中
\begin{equation}
  F_{\alpha_1,\alpha_2,Q/2+\ii p,\alpha_3,\alpha_4}(q)
  = \sum_{\{n\},\{n'\}} q^{\sum n}\,
  l_{\alpha_1,\alpha_2,Q/2+\ii p,\{n\}}\,
  G^{\{n\},\{n'\}}_p\,
  r_{\{\tilde n'\},Q/2-\ii p,\alpha_3,\alpha_4}
  \tag{2.60}
\end{equation}
称为四点共形块。

% [Tutor Note: p.14, 式 (2.60) 的指标：求和指标是 $\{n\},\{n'\}$，但最后一个因子 $r$ 的下标在原文
%   印作 $\{\tilde n'\}$（与 (2.58) 中 $r_{\{\tilde n'\},\cdots}$ 的写法一致）。按 (2.55)→(2.59) 的缩并
%   方式，这里参与缩并的应当是 $\{n'\}$：$l$ 带 $\{n\}$、$G$ 带 $\{n\},\{n'\}$、$r$ 带 $\{n'\}$。
%   此处已用 800 dpi 裁图复核，原文确实印作 $\{\tilde n'\}$，疑为 $\{n'\}$ 之误；正文照原样排印，
%   读者按 $r_{\{n'\},\cdots}$ 理解即可。]

% [Tutor Note: p.14, 式 (2.55) 中 $z^{-\sum n}\bar z^{-\sum\tilde n}|z|^{-2\Delta_p}$ 三项的分工：
%   $|z|^{-2\Delta_p}$ 是主场（$p$ 标记的中间态）自身的维数因子，$z^{-\sum n}\bar z^{-\sum\tilde n}$
%   是每个后裔 $L_{-n}$、$\bar L_{-\tilde n}$ 贡献的坐标幂。合起来使总能级
%   $N=\sum n+\sum\tilde n$ 的项带 $q^{N}$（见 (2.60) 的 $q^{\sum n}$），所以 (2.72) 的一阶项只需
%   用到 $l_{\cdots,\{1\}}$、$G^{\{1\},\{1\}}$、$r_{\{1\},\cdots}$ 这一组数据。]

注意，表达式 (2.59) 是在 $z=\infty,1$ 处的配对 $\alpha_1,\alpha_2$ 与 $z=q,0$ 处的配对
$\alpha_3,\alpha_4$ 之间插入完备态集得到的。同一个关联函数也可以通过在下述配对之间插入完备态集
得到：$z=\infty,q$ 处的配对 $\alpha_1,\alpha_3$ 与 $z=1,0$ 处的配对 $\alpha_2,\alpha_4$；
或者 $z=\infty,0$ 处的配对 $\alpha_1,\alpha_4$ 与 $z=1,q$ 处的配对 $\alpha_2,\alpha_3$。
所得表达式之间的相等绝非平凡，但已在 [16,19] 中得到证明。

% [Tutor Note: p.14, 末段：这就是共形场论的交叉对称性（crossing symmetry），也就是 OPE 的结合律。
%   三种配对对应 (2.59) 的三种"通道"展开，它们必须给出同一个关联函数；(2.60) 的共形块在不同通道
%   之间通过 crossing kernel 相联系。原文这里只声明"已在 [16,19] 中证明"，未给证明，属综述性引用，
%   不必在此补推导；但要注意 (2.59) 的三个表达式都含 $\int_0^\infty\dd p$ 与连续谱，所以"相等"是
%   逐点（对 $q,\bar q$）成立的恒等式，而不是逐项相消。]
